\chapter{venv 环境}

\section{概述}
\codeinline{venv}就是用来为一个应用创建一套“隔离”的Python运行环境

\section{使用}
\begin{lstlisting}[style=demostyle,numbers=none]
# 创建虚拟环境
python -m venv my_env

# 激活虚拟环境
## Linux
source my_env/bin/activate
## Windows
my_env\Scripts\activate.bat

# 退出虚拟环境
deactivate
\end{lstlisting}

\section{常用命令}
\subsection{指定python版本}
\begin{lstlisting}[style=demostyle,numbers=none]
# 创建虚拟环境
python -m venv my_env --python=python3.6
\end{lstlisting}

\section{与Vscode结合}
\url{https://code.visualstudio.com/docs/python/environments}
